\documentclass[11pt,a4paper]{article}

\usepackage[utf8]{inputenc}
\usepackage[T1]{fontenc}
\usepackage[french]{babel}
\usepackage[a4paper,margin=2.2cm]{geometry}
\usepackage{amsmath,amssymb,amsthm,mathtools}
\usepackage{lmodern}
\usepackage{enumitem}
\usepackage{hyperref}
\usepackage{xcolor}
\usepackage{fancyhdr}
\usepackage{stmaryrd}
\usepackage{mathrsfs}
\usepackage{array}

\hypersetup{
    colorlinks=true,
    linkcolor=black,
    urlcolor=blue,
    citecolor=black
}

\setlength{\parindent}{0pt}
\setlength{\parskip}{0.6em}

\pagestyle{fancy}
\fancyhf{}
\fancyhead[L]{Géométrie --- Chapitre 13}
\fancyhead[R]{Géométrie affine et euclidienne du plan et de l’espace}
\fancyfoot[C]{\thepage}

% Environnements
\newtheorem{definition}{Définition}[section]
\newtheorem{theoreme}[definition]{Théorème}
\newtheorem{proposition}[definition]{Proposition}
\newtheorem{corollaire}[definition]{Corollaire}
\newtheorem{lemme}[definition]{Lemme}
\theoremstyle{remark}
\newtheorem{remarque}[definition]{Remarque}
\newtheorem{exemple}[definition]{Exemple}
\newtheorem*{methode}{Méthode}
\newtheorem*{idee}{Idée}

% Commandes
\newcommand{\R}{\mathbb{R}}
\newcommand{\Vect}{\mathrm{Vect}}
\newcommand{\scal}[2]{\left\langle #1,#2\right\rangle}
\newcommand{\Norm}[1]{\left\|#1\right\|}
\newcommand{\abs}[1]{\left|#1\right|}
\newcommand{\uu}{\overrightarrow{u}}
\newcommand{\vv}{\overrightarrow{v}}
\newcommand{\ww}{\overrightarrow{w}}

\begin{document}

\begin{center}
    {\LARGE \textbf{Chapitre 13 --- Géométrie affine et euclidienne du plan et de l’espace}}\\[0.4em]
    {\large Droites, plans, équations, distances et interprétation vectorielle}
\end{center}

\hrule
\vspace{1em}

\tableofcontents

\newpage

\section{Pourquoi ce chapitre ?}

Après l’algèbre linéaire, il est naturel de revenir à la géométrie concrète :
\begin{itemize}
    \item décrire une droite ou un plan ;
    \item calculer une intersection ;
    \item reconnaître des positions relatives ;
    \item mesurer une distance ou un angle.
\end{itemize}

L’intérêt de ce chapitre est double :
\begin{itemize}
    \item donner du sens géométrique aux notions vectorielles ;
    \item fournir des outils de calcul très efficaces.
\end{itemize}

\begin{idee}
La géométrie affine décrit les \emph{positions} et les \emph{ensembles}.  
La géométrie euclidienne ajoute les \emph{longueurs}, les \emph{angles} et les \emph{distances}.
\end{idee}

\section{Cadre affine}

\subsection{Points et vecteurs}

\begin{remarque}
Dans le plan ou l’espace, il faut distinguer :
\begin{itemize}
    \item les \textbf{points}, qui représentent des positions ;
    \item les \textbf{vecteurs}, qui représentent des déplacements.
\end{itemize}
\end{remarque}

\begin{definition}
Si \(A\) et \(B\) sont deux points, on note
\[
\overrightarrow{AB}
\]
le vecteur allant de \(A\) vers \(B\).
\end{definition}

\begin{proposition}
Pour tous points \(A,B,C\),
\[
\overrightarrow{AB}+\overrightarrow{BC}=\overrightarrow{AC}.
\]
\end{proposition}

\begin{proof}
C’est la relation de Chasles.
\end{proof}

\begin{remarque}
Toute la géométrie affine peut être reformulée à l’aide de relations vectorielles.
\end{remarque}

\subsection{Repères}

\begin{definition}
Dans le plan affine, un \textbf{repère} est la donnée d’un point \(O\) et d’une base \((\vec{i},\vec{j})\) de \(\R^2\).

Dans l’espace affine, un \textbf{repère} est la donnée d’un point \(O\) et d’une base \((\vec{i},\vec{j},\vec{k})\) de \(\R^3\).
\end{definition}

\begin{remarque}
Dans un repère, tout point \(M\) est repéré par les coordonnées du vecteur \(\overrightarrow{OM}\).
\end{remarque}

\section{Barycentre et milieux}

\subsection{Milieu}

\begin{definition}
Le \textbf{milieu} du segment \([AB]\) est le point \(I\) tel que
\[
\overrightarrow{OI}=\frac12\big(\overrightarrow{OA}+\overrightarrow{OB}\big).
\]
\end{definition}

\begin{proposition}
Le point \(I\) est milieu de \([AB]\) si et seulement si
\[
\overrightarrow{IA}+\overrightarrow{IB}=0.
\]
\end{proposition}

\subsection{Barycentre de deux points}

\begin{definition}
Soient \(A,B\) deux points et \(\alpha,\beta\in\R\) tels que
\[
\alpha+\beta\neq 0.
\]
Le barycentre de \((A,\alpha)\) et \((B,\beta)\) est le point \(G\) défini par
\[
\overrightarrow{OG}=\frac{\alpha\overrightarrow{OA}+\beta\overrightarrow{OB}}{\alpha+\beta}.
\]
\end{definition}

\begin{remarque}
Le milieu est le cas particulier \(\alpha=\beta=1\).
\end{remarque}

\section{Droites affines du plan et de l’espace}

\subsection{Définition vectorielle}

\begin{definition}
Soit \(A\) un point et \(\vec{u}\neq 0\) un vecteur.  
La \textbf{droite affine} passant par \(A\) et dirigée par \(\vec{u}\) est l’ensemble
\[
\mathcal D=A+\Vect(\vec{u})
=
\{M\mid \overrightarrow{AM}=t\vec{u},\ t\in\R\}.
\]
\end{definition}

\begin{remarque}
Une droite est donc un point de départ plus une direction.
\end{remarque}

\subsection{Représentation paramétrique}

\begin{proposition}
Si \(A(x_A,y_A)\) dans le plan et \(\vec{u}=(a,b)\), alors une représentation paramétrique de la droite est
\[
\begin{cases}
x=x_A+ta,\\
y=y_A+tb,
\end{cases}
\qquad t\in\R.
\]
\end{proposition}

Dans l’espace, si \(A(x_A,y_A,z_A)\) et \(\vec{u}=(a,b,c)\), on a
\[
\begin{cases}
x=x_A+ta,\\
y=y_A+tb,\\
z=z_A+tc,
\end{cases}
\qquad t\in\R.
\]

\begin{exemple}
La droite passant par \(A(1,2)\) et de vecteur directeur \((3,-1)\) admet pour représentation
\[
\begin{cases}
x=1+3t,\\
y=2-t.
\end{cases}
\]
\end{exemple}

\subsection{Équation cartésienne dans le plan}

\begin{theoreme}
Toute droite du plan admet une équation cartésienne de la forme
\[
ax+by+c=0,
\]
où \((a,b)\neq (0,0)\).
\end{theoreme}

\begin{proof}
Soit une droite \(\mathcal D\) passant par \(A(x_A,y_A)\) et dirigée par \((u,v)\neq (0,0)\).  
Un vecteur normal à la droite est \((a,b)=(-v,u)\). Alors
\[
M(x,y)\in\mathcal D
\Longleftrightarrow
\scal{(a,b)}{\overrightarrow{AM}}=0.
\]
Cela donne
\[
a(x-x_A)+b(y-y_A)=0,
\]
soit
\[
ax+by+c=0.
\]
\end{proof}

\begin{remarque}
Le vecteur \((a,b)\) est un \textbf{vecteur normal} à la droite.
\end{remarque}

\begin{exemple}
La droite d’équation
\[
2x-3y+1=0
\]
a pour vecteur normal \((2,-3)\), donc un vecteur directeur peut être \((3,2)\).
\end{exemple}

\section{Plans affines de l’espace}

\subsection{Définition vectorielle}

\begin{definition}
Soit \(A\) un point et \(\vec{u},\vec{v}\) deux vecteurs non colinéaires.  
Le \textbf{plan affine} passant par \(A\) et dirigé par \(\vec{u},\vec{v}\) est
\[
\mathcal P=A+\Vect(\vec{u},\vec{v})
=
\{M\mid \overrightarrow{AM}=s\vec{u}+t\vec{v},\ s,t\in\R\}.
\]
\end{definition}

\subsection{Représentation paramétrique}

\begin{proposition}
Si \(A(x_A,y_A,z_A)\), \(\vec{u}=(a,b,c)\), \(\vec{v}=(a',b',c')\), alors une représentation paramétrique du plan est
\[
\begin{cases}
x=x_A+sa+ta',\\
y=y_A+sb+tb',\\
z=z_A+sc+tc',
\end{cases}
\qquad (s,t)\in\R^2.
\]
\end{proposition}

\subsection{Équation cartésienne d’un plan}

\begin{theoreme}
Tout plan affine de l’espace admet une équation cartésienne de la forme
\[
ax+by+cz+d=0,
\]
où \((a,b,c)\neq (0,0,0)\).
\end{theoreme}

\begin{proof}
Le vecteur \((a,b,c)\) est un vecteur normal au plan.  
Si \(A(x_A,y_A,z_A)\in\mathcal P\), alors
\[
M(x,y,z)\in\mathcal P
\Longleftrightarrow
a(x-x_A)+b(y-y_A)+c(z-z_A)=0.
\]
On obtient alors
\[
ax+by+cz+d=0.
\]
\end{proof}

\begin{remarque}
Le vecteur \((a,b,c)\) est un \textbf{vecteur normal} au plan.
\end{remarque}

\section{Positions relatives}

\subsection{Deux droites du plan}

\begin{proposition}
Deux droites du plan sont :
\begin{itemize}
    \item soit sécantes ;
    \item soit parallèles distinctes ;
    \item soit confondues.
\end{itemize}
\end{proposition}

\begin{remarque}
Deux droites du plan sont parallèles si et seulement si leurs vecteurs directeurs sont colinéaires.
\end{remarque}

\subsection{Deux droites de l’espace}

\begin{proposition}
Dans l’espace, deux droites peuvent être :
\begin{itemize}
    \item sécantes ;
    \item parallèles ;
    \item confondues ;
    \item gauches.
\end{itemize}
\end{proposition}

\begin{definition}
Deux droites sont dites \textbf{gauches} si elles ne sont ni parallèles ni sécantes.
\end{definition}

\subsection{Droite et plan}

\begin{proposition}
Une droite et un plan de l’espace peuvent être :
\begin{itemize}
    \item sécants en un point ;
    \item parallèles ;
    \item la droite peut être contenue dans le plan.
\end{itemize}
\end{proposition}

\subsection{Deux plans}

\begin{proposition}
Deux plans de l’espace peuvent être :
\begin{itemize}
    \item sécants selon une droite ;
    \item parallèles distincts ;
    \item confondus.
\end{itemize}
\end{proposition}

\section{Intersections}

\subsection{Intersection de deux droites du plan}

\begin{methode}
Pour déterminer l’intersection de deux droites du plan :
\begin{enumerate}
    \item écrire leurs équations cartésiennes ;
    \item résoudre le système de deux équations à deux inconnues ;
    \item conclure : une solution, aucune, ou infinité.
\end{enumerate}
\end{methode}

\begin{exemple}
Déterminons l’intersection de
\[
\begin{cases}
x+y=3,\\
2x-y=0.
\end{cases}
\]
On obtient
\[
3x=3 \Longrightarrow x=1,
\qquad y=2.
\]
Les droites sont donc sécantes en \((1,2)\).
\end{exemple}

\subsection{Intersection droite-plan}

\begin{methode}
Pour étudier l’intersection d’une droite et d’un plan :
\begin{enumerate}
    \item écrire une représentation paramétrique de la droite ;
    \item remplacer dans l’équation du plan ;
    \item résoudre en le paramètre.
\end{enumerate}
\end{methode}

\begin{exemple}
Soit la droite
\[
\begin{cases}
x=1+t,\\
y=2-t,\\
z=t
\end{cases}
\]
et le plan
\[
x+y+z-4=0.
\]
En remplaçant :
\[
(1+t)+(2-t)+t-4=0
\Longrightarrow
t-1=0
\Longrightarrow
t=1.
\]
Le point d’intersection est
\[
(2,1,1).
\]
\end{exemple}

\section{Produit scalaire en géométrie}

\subsection{Rappel}

Dans \(\R^2\) ou \(\R^3\), avec le produit scalaire canonique,
\[
\scal{u}{v}=x_1y_1+x_2y_2
\quad \text{ou} \quad
x_1y_1+x_2y_2+x_3y_3.
\]

\subsection{Orthogonalité}

\begin{proposition}
Deux vecteurs non nuls \(u\) et \(v\) sont orthogonaux si et seulement si
\[
\scal{u}{v}=0.
\]
\end{proposition}

\subsection{Angle entre deux vecteurs}

\begin{definition}
Si \(u\) et \(v\) sont deux vecteurs non nuls, on définit l’angle \(\theta\in[0,\pi]\) par
\[
\cos\theta=\frac{\scal{u}{v}}{\Norm{u}\Norm{v}}.
\]
\end{definition}

\begin{remarque}
Cette formule est valable grâce à l’inégalité de Cauchy-Schwarz.
\end{remarque}

\begin{exemple}
Pour
\[
u=(1,1),\qquad v=(1,-1),
\]
on a
\[
\scal{u}{v}=1-1=0,
\]
donc l’angle vaut \(\frac{\pi}{2}\).
\end{exemple}

\subsection{Vecteur normal et produit scalaire}

\begin{proposition}
Dans le plan, la droite d’équation
\[
ax+by+c=0
\]
admet pour vecteur normal \((a,b)\).

Dans l’espace, le plan d’équation
\[
ax+by+cz+d=0
\]
admet pour vecteur normal \((a,b,c)\).
\end{proposition}

\begin{proof}
Cela découle directement des formes
\[
a(x-x_0)+b(y-y_0)=0
\]
et
\[
a(x-x_0)+b(y-y_0)+c(z-z_0)=0.
\]
\end{proof}

\section{Distances}

\subsection{Distance entre deux points}

\begin{definition}
Si \(A\) et \(B\) sont deux points, la distance \(AB\) est
\[
AB=\Norm{\overrightarrow{AB}}.
\]
\end{definition}

Dans \(\R^2\),
\[
AB=\sqrt{(x_B-x_A)^2+(y_B-y_A)^2}.
\]

Dans \(\R^3\),
\[
AB=\sqrt{(x_B-x_A)^2+(y_B-y_A)^2+(z_B-z_A)^2}.
\]

\subsection{Distance d’un point à une droite du plan}

\begin{theoreme}
Soit \(M_0(x_0,y_0)\) un point et \(\mathcal D\) la droite d’équation
\[
ax+by+c=0,
\qquad (a,b)\neq (0,0).
\]
Alors
\[
d(M_0,\mathcal D)=\frac{\abs{ax_0+by_0+c}}{\sqrt{a^2+b^2}}.
\]
\end{theoreme}

\begin{proof}
Le vecteur \((a,b)\) est normal à la droite. La distance est la valeur absolue de la projection du vecteur reliant un point de la droite à \(M_0\) sur ce vecteur normal unitaire. On obtient exactement la formule ci-dessus.
\end{proof}

\begin{exemple}
Distance du point \(M(1,2)\) à la droite
\[
2x-y+3=0 :
\]
\[
d(M,\mathcal D)=\frac{\abs{2\cdot 1-2+3}}{\sqrt{2^2+(-1)^2}}
=\frac{3}{\sqrt5}.
\]
\end{exemple}

\subsection{Distance d’un point à un plan}

\begin{theoreme}
Soit \(M_0(x_0,y_0,z_0)\) et le plan
\[
\mathcal P: ax+by+cz+d=0,
\qquad (a,b,c)\neq (0,0,0).
\]
Alors
\[
d(M_0,\mathcal P)=\frac{\abs{ax_0+by_0+cz_0+d}}{\sqrt{a^2+b^2+c^2}}.
\]
\end{theoreme}

\begin{proof}
Même idée que dans le plan : on projette selon un vecteur normal au plan.
\end{proof}

\section{Produit vectoriel dans \(\R^3\)}

\subsection{Définition}

\begin{definition}
Pour
\[
u=(x,y,z),\qquad v=(x',y',z'),
\]
on définit le \textbf{produit vectoriel}
\[
u\wedge v=
\begin{pmatrix}
yz'-zy'\\
zx'-xz'\\
xy'-yx'
\end{pmatrix}.
\]
\end{definition}

\subsection{Propriétés}

\begin{proposition}
Le vecteur \(u\wedge v\) est orthogonal à \(u\) et à \(v\).
\end{proposition}

\begin{proof}
Un calcul direct montre que
\[
\scal{u\wedge v}{u}=0
\qquad \text{et} \qquad
\scal{u\wedge v}{v}=0.
\]
\end{proof}

\begin{proposition}
On a
\[
u\wedge v=0
\Longleftrightarrow
u \text{ et } v \text{ sont colinéaires.}
\]
\end{proposition}

\begin{proof}
Si \(u\) et \(v\) sont colinéaires, le produit vectoriel est nul par bilinéarité et antisymétrie.

Réciproquement, si \(u\wedge v=0\), alors \(u\) et \(v\) sont dépendants dans \(\R^3\), donc colinéaires.
\end{proof}

\subsection{Norme du produit vectoriel}

\begin{proposition}
Si \(\theta\) est l’angle entre \(u\) et \(v\), alors
\[
\Norm{u\wedge v}=\Norm{u}\Norm{v}\sin\theta.
\]
\end{proposition}

\begin{remarque}
La norme du produit vectoriel mesure l’aire du parallélogramme engendré par \(u\) et \(v\).
\end{remarque}

\subsection{Équation d’un plan via un produit vectoriel}

\begin{proposition}
Si un plan est dirigé par deux vecteurs non colinéaires \(u\) et \(v\), alors
\[
u\wedge v
\]
est un vecteur normal au plan.
\end{proposition}

\begin{proof}
Le produit vectoriel est orthogonal à \(u\) et à \(v\), donc à tout vecteur du plan.
\end{proof}

\section{Équations paramétriques et cartésiennes : passage de l’une à l’autre}

\subsection{Droite du plan}

\begin{methode}
Pour passer d’une représentation paramétrique
\[
\begin{cases}
x=x_0+at,\\
y=y_0+bt
\end{cases}
\]
à une équation cartésienne :
\begin{enumerate}
    \item éliminer le paramètre \(t\) ;
    \item ou utiliser un vecteur normal \((-b,a)\).
\end{enumerate}
\end{methode}

\subsection{Plan de l’espace}

\begin{methode}
Pour obtenir une équation cartésienne d’un plan à partir de deux vecteurs directeurs \(u,v\) :
\begin{enumerate}
    \item calculer un vecteur normal \(n=u\wedge v\) ;
    \item écrire
    \[
    \scal{n}{\overrightarrow{AM}}=0.
    \]
\end{enumerate}
\end{methode}

\section{Positions relatives via les vecteurs normaux}

\subsection{Deux droites du plan}

\begin{proposition}
Deux droites du plan sont parallèles si et seulement si leurs vecteurs normaux sont colinéaires.
\end{proposition}

\subsection{Deux plans de l’espace}

\begin{proposition}
Deux plans de l’espace sont parallèles si et seulement si leurs vecteurs normaux sont colinéaires.
\end{proposition}

\subsection{Droite et plan}

\begin{proposition}
Une droite de vecteur directeur \(u\) est parallèle à un plan de vecteur normal \(n\) si et seulement si
\[
\scal{u}{n}=0.
\]
\end{proposition}

\begin{proof}
La droite est parallèle au plan si son vecteur directeur appartient à la direction du plan, donc est orthogonal au vecteur normal du plan.
\end{proof}

\section{Systèmes cartésiens et interprétation géométrique}

\subsection{Dans le plan}

Un système de deux équations cartésiennes à deux inconnues représente l’intersection de deux droites.

\subsection{Dans l’espace}

\begin{itemize}
    \item une équation cartésienne représente un plan ;
    \item deux équations représentent l’intersection de deux plans, souvent une droite ;
    \item trois équations représentent l’intersection de trois plans.
\end{itemize}

\begin{remarque}
L’étude géométrique des positions relatives se ramène donc souvent à la résolution de systèmes linéaires.
\end{remarque}

\section{Interprétation vectorielle de la colinéarité et de l’alignement}

\subsection{Colinéarité}

\begin{definition}
Deux vecteurs \(u\) et \(v\) sont colinéaires s’il existe \(\lambda\in\R\) tel que
\[
u=\lambda v.
\]
\end{definition}

\subsection{Alignement}

\begin{proposition}
Trois points \(A,B,C\) sont alignés si et seulement si
\[
\overrightarrow{AB}
\quad \text{et} \quad
\overrightarrow{AC}
\]
sont colinéaires.
\end{proposition}

\subsection{Coplanarité dans l’espace}

\begin{proposition}
Quatre points \(A,B,C,D\) sont coplanaires si et seulement si les vecteurs
\[
\overrightarrow{AB},\quad \overrightarrow{AC},\quad \overrightarrow{AD}
\]
sont liés.
\end{proposition}

\begin{proof}
Les points sont coplanaires si et seulement si \(D\) appartient au plan affine passant par \(A\) dirigé par \(\overrightarrow{AB}\) et \(\overrightarrow{AC}\), c’est-à-dire si
\[
\overrightarrow{AD}\in \Vect(\overrightarrow{AB},\overrightarrow{AC}).
\]
\end{proof}

\section{Angles géométriques}

\subsection{Angle de deux droites}

\begin{definition}
L’angle de deux droites est l’angle de deux de leurs vecteurs directeurs non nuls.
\end{definition}

\subsection{Angle de deux plans}

\begin{definition}
L’angle de deux plans est l’angle de deux vecteurs normaux non nuls.
\end{definition}

\subsection{Angle d’une droite et d’un plan}

\begin{definition}
L’angle \(\theta\) entre une droite de vecteur directeur \(u\) et un plan de vecteur normal \(n\) est défini par
\[
\sin\theta=\frac{\abs{\scal{u}{n}}}{\Norm{u}\Norm{n}}.
\]
\end{definition}

\begin{remarque}
Quand \(u\) est orthogonal à \(n\), la droite est parallèle au plan et \(\theta=0\).
\end{remarque}

\section{Méthodes pratiques}

\begin{methode}
Pour étudier une droite ou un plan :
\begin{enumerate}
    \item identifier un point ;
    \item identifier des vecteurs directeurs ou un vecteur normal ;
    \item choisir la forme la plus adaptée : paramétrique ou cartésienne.
\end{enumerate}
\end{methode}

\begin{methode}
Pour calculer une distance :
\begin{enumerate}
    \item repérer un vecteur normal si l’objet est donné par une équation cartésienne ;
    \item utiliser la formule de projection sur ce vecteur normal.
\end{enumerate}
\end{methode}

\begin{methode}
Pour étudier une intersection :
\begin{enumerate}
    \item remplacer une représentation paramétrique dans une équation cartésienne ;
    \item ou résoudre directement le système obtenu.
\end{enumerate}
\end{methode}

\section{Erreurs classiques}

\begin{itemize}
    \item Confondre vecteur directeur et vecteur normal.
    \item Oublier qu’une droite dans l’espace n’admet pas en général une seule équation cartésienne.
    \item Se tromper entre parallélisme et orthogonalité.
    \item Utiliser la formule de distance sans vérifier que les coefficients du vecteur normal ne sont pas tous nuls.
    \item Confondre produit scalaire nul et égalité de vecteurs.
    \item Oublier que deux droites de l’espace peuvent être gauches.
\end{itemize}

\section{Résumé du chapitre}

\begin{itemize}
    \item Une droite affine s’écrit
    \[
    A+\Vect(u).
    \]
    \item Un plan affine s’écrit
    \[
    A+\Vect(u,v).
    \]
    \item Dans le plan, une droite admet une équation cartésienne
    \[
    ax+by+c=0.
    \]
    \item Dans l’espace, un plan admet une équation cartésienne
    \[
    ax+by+cz+d=0.
    \]
    \item Le produit scalaire permet de caractériser l’orthogonalité et les angles.
    \item La distance d’un point à une droite ou à un plan s’exprime à l’aide d’un vecteur normal.
    \item Dans \(\R^3\), le produit vectoriel fournit un vecteur normal à deux vecteurs non colinéaires.
    \item Les problèmes de géométrie affine se ramènent souvent à des systèmes linéaires.
\end{itemize}

\section*{Exercices d’application immédiate}

\begin{enumerate}[label=\textbf{Exercice \arabic*.}]
    \item Donner une représentation paramétrique de la droite passant par \(A(1,2)\) et dirigée par \((3,-1)\).

    \item Déterminer une équation cartésienne de la droite passant par \(A(1,2)\) et \(B(3,5)\).

    \item Étudier la position relative des droites
    \[
    2x-y+1=0
    \qquad \text{et} \qquad
    4x-2y+3=0.
    \]

    \item Donner une représentation paramétrique du plan passant par \(A(1,0,2)\) et dirigé par
    \[
    u=(1,1,0),\qquad v=(0,1,1).
    \]

    \item Déterminer une équation cartésienne du plan passant par \(A(1,0,2)\) et dirigé par les vecteurs précédents.

    \item Déterminer l’intersection de la droite
    \[
    \begin{cases}
    x=1+t,\\
    y=2-t,\\
    z=t
    \end{cases}
    \]
    avec le plan
    \[
    x+y+z-4=0.
    \]

    \item Calculer la distance du point \(M(1,2)\) à la droite
    \[
    2x-y+3=0.
    \]

    \item Calculer la distance du point \(M(1,2,3)\) au plan
    \[
    x-y+2z-4=0.
    \]

    \item Calculer le produit vectoriel de
    \[
    u=(1,2,3),\qquad v=(0,1,-1).
    \]

    \item Montrer que quatre points \(A,B,C,D\) de l’espace sont coplanaires si et seulement si
    \[
    \overrightarrow{AD}\in \Vect(\overrightarrow{AB},\overrightarrow{AC}).
    \]
\end{enumerate}

\end{document}